\begin{frame}{Pensiones}
\begin{itemize}
%    \item El gasto pensional no se limita a Colpensiones: incluye regímenes especiales (FFAA, docentes, empleados públicos).
%    \item Las asignaciones de retiro de FFAA siguen vigentes y presionan las finanzas públicas.
%    \item 2024–2035: gasto en pensiones +1 pp PIB (Colpensiones +0,7 pp; asignaciones de retiro +0,3 pp).
    \item Baja cobertura: sólo 1 de cada 4 adultos mayores recibe pensión.
%\end{itemize}
%\end{frame}

%\begin{frame}{Pensiones}
%\begin{enumerate}
    \item Ajustes paramétricos: 
    \begin{itemize}
        \item Edad de pensión
        \item Tasa de reemplazo
        \item Umbral de cotización
    \end{itemize}    
    %\item Contener el aumento del salario mínimo.
    \item Eliminar la exención de impuestos a las pensiones.
    \item Aumentar cotización en salud de los pensionados.
    \item Incrementar aportes de los pensionados al Fondo de Solidaridad Pensional.
    %\item Eliminar pensiones múltiples financiadas con recursos públicos.
\end{itemize}
\end{frame}

%\begin{frame}{Pensiones}
%\begin{itemize}
%    \item Ley 2381 de 2024: más cotizaciones al componente público; fondo de ahorro en BanRep.
%    \item Riesgo de uso de los recursos del fondo ante presiones fiscales → agotamiento proyectado en 2064.
%    \item Necesarios ajustes paramétricos: umbral de ingresos, tasa de reemplazo, edad de jubilación por cohortes.
%    \item Implementar pronto → señal de compromiso fiscal; reducción de pasivo implícito y riesgo país.
%\end{itemize}
%\end{frame}
